% !TEX program = pdflatex
\documentclass[11pt,a4paper]{article}

\usepackage[T1]{fontenc}
\usepackage[utf8]{inputenc}
\usepackage{lmodern}
\usepackage{amsmath,amssymb,amsthm}
\usepackage{geometry}
\geometry{margin=2.6cm}
\usepackage{booktabs}
\usepackage{xcolor}
\usepackage[colorlinks=true,linkcolor=blue!55!black,citecolor=blue!55!black,
            urlcolor=blue!55!black]{hyperref}
\usepackage{enumitem}
\usepackage{mdframed}

\linespread{1.10}
\setlength{\emergencystretch}{2em}

\newtheorem{theorem}{Theorem}[section]
\newtheorem{proposition}[theorem]{Proposition}
\newtheorem{corollary}[theorem]{Corollary}
\theoremstyle{definition}
\newtheorem{definition}[theorem]{Definition}
\theoremstyle{remark}
\newtheorem{remark}[theorem]{Remark}

\newcommand{\lean}[1]{\texttt{\small #1}}
\newcommand{\Ric}{\mathrm{Ric}}
\newcommand{\Rsc}{\mathrm{R}}
\newcommand{\Defm}{D}
\newcommand{\sg}{\sigma}
\newcommand{\Ap}{\ensuremath{\mathrm{A6}'}}
\newcommand{\App}{\ensuremath{\mathrm{A6}''}}
\newcommand{\Afour}{\ensuremath{\mathrm{A4}'}}

\newmdenv[linewidth=0.6pt,linecolor=black!45,backgroundcolor=black!3,
          skipabove=8pt,skipbelow=8pt,innertopmargin=6pt,innerbottommargin=6pt]{claimbox}

\title{\bf Scale-Coupled Dynamics II:\\
the pairing, the source law, and a torsion-free antisymmetric Ricci}

\author{}
\date{}

\begin{document}
\maketitle

\begin{abstract}
Scale-Coupled Dynamics (SCD) is a formalised framework in which the only
primitive is a local scale and the geometry is its deformation. This note
reports a revision pass carried out entirely inside the Lean~4 formalisation.
Its results are of three kinds.

\emph{Structural.} The founding statement \emph{scale $\times$ energy $=$
constant} needs only a group homomorphism; the commutative ring the development
was built on was silently supplying a second, independent thing --- a bilinear
pairing for the geometry. Separating them consolidates the axioms from seven to
five, and shows that the pairing was \emph{never chosen}: the gravitational
coupling was computed with a Euclidean trace while the observables live in a
Lorentzian sector, joined by a single theorem whose hypotheses are mutually
exclusive.

\emph{Corrective.} Redoing the curvature identity with a signature shows the
coefficient $\kappa=-2(n-1)\Delta$ survives unchanged, for a structural reason,
but the operator it multiplies is the d'Alembertian rather than the Laplacian:
the source law is hyperbolic, and Poisson's equation is its static limit. Two
further claims lose their stated grounds while keeping their conclusions, and one
claim is withdrawn.

\emph{Predictive.} Off the commuting locus the Ricci tensor is \emph{not
symmetric}, with $\Ric_{be}-\Ric_{eb}=(n-2)[\sg_b,\sg_e]$, while the connection
remains torsion-free and the scalar curvature does not move at all. Nothing
classical --- Newton's law, the coupling, cosmology --- takes a quantum
correction at any order; in particular there is no $\hbar G/r^{3}$ correction to
the Newtonian potential. Finally, the usual four-constants--three-dimensions
bookkeeping collapses to a single grading in which $\hbar$ and $G$ carry the
\emph{same} weight, so that ``are they the same magnitude?'' becomes a
well-formed question whose answer is one pure number --- and the framework's
long-quoted ``one free pure number'' is conditional on answering it.

All statements below are theorems in a Lean~4 development of 1275 theorems, every
one of which passes \lean{\#print axioms} with no axiom beyond \lean{propext},
\lean{Classical.choice} and \lean{Quot.sound}, and with no \lean{sorry}.
\end{abstract}

\tableofcontents

\section{Scope, and how to read this}

SCD posits a dimensionless local log-scale $\sg$ over a ring with $n$ commuting
derivations, and derives geometry as the deformation of that scale rather than as
an independent metric structure. This note does not re-derive the framework. It
reports what a deliberate examination of its \emph{mathematical form} produced,
and it is written so that each claim can be checked against a named Lean
declaration.

Three conventions are used throughout, and they are not decoration.

\begin{itemize}[leftmargin=1.4em]
\item \textbf{Claimed} means proved in the formalisation, with the declaration
      named.
\item \textbf{Identified} means an interpretive step that the formalisation
      records as an input rather than a consequence.
\item \textbf{Withdrawn} means a claim previously made in this development and
      now retracted, with the reason given.
\end{itemize}

Section~\ref{sec:open} collects what remains open, including one item this pass
deliberately declined to pursue.

\section{The founding statement needs a group, not a ring}
\label{sec:reduction}

The framework's founding statement is \emph{scale $\times$ energy $=$ constant},
encoded by two axioms: A2 supplies a log-scale $\sg$ with $s=e^{\sg}$ a unit, and
A3 sets $\varepsilon = s^{-1}$.

The theorem saying what this means had never been written.

\begin{theorem}[\lean{Reduction.logDeriv\_mul}]
For units $s,t$ of the ring and each direction $i$,
\[
  \operatorname{logDeriv}(st)_i \;=\; \operatorname{logDeriv}(s)_i +
  \operatorname{logDeriv}(t)_i .
\]
\end{theorem}

The log-derivative $s^{-1}\partial_i s$ is a \emph{group homomorphism} from the
units into the additive group of gradients. Two consequences follow at once.

\begin{corollary}[\lean{Reduction.logDeriv\_inv}]
$\operatorname{logDeriv}(s^{-1}) = -\operatorname{logDeriv}(s)$.
\end{corollary}

This \emph{is} A3. That the energy's gradient is minus the scale's is inversion in
a group, not an independent postulate.

\begin{definition}[\lean{Reduction.ScaleHom}]
A group $G$, an abelian group $M$, and a map $D\colon G\to M^{n}$ with
$D(st)=Ds+Dt$.
\end{definition}

Every scale algebra furnishes one (\lean{Reduction.ofScaleAlgebra}). The signature
contains \emph{no ring, no commutativity, and no multiplication}. That is the
whole of what the founding statement requires.

\subsection{So what was the ring doing?}

\begin{theorem}[\lean{Reduction.sig\_add}, \lean{Reduction.gradsq\_add}]
The gradient is additive in the log-scale, while
\[
  \operatorname{gradsq}(\sg+\tau)
   = \operatorname{gradsq}\sg + \operatorname{gradsq}\tau
     + 2\sum_i \operatorname{sig}\sg_i\cdot\operatorname{sig}\tau_i .
\]
\end{theorem}

The axiom is linear in the gradient; the geometry is quadratic; and the failure of
additivity is exactly a \emph{symmetric bilinear pairing}. Every geometric object
in the development is built from that pairing.

\begin{claimbox}
A commutative ring supplies the homomorphism and the pairing in one move. That is
why it was chosen, and it was never recorded that these are two commitments rather
than one.
\end{claimbox}

The sharp form of the concern: in a ring, the multiplication defining the unit
group and the multiplication supplying the pairing are the \emph{same operation},
and nothing in \emph{scale $\times$ energy $=$ constant} requires that.

\section{Axiom consolidation: seven to five}
\label{sec:consolidation}

Three results turn a long-standing goal into a theorem.

\begin{proposition}[\lean{Consolidation.A3\_is\_inversion}]
The duality relation is \lean{Units.inv\_mul} and $\varepsilon:=s^{-1}$ is a
definition. A3 names a definition and asserts a group law.
\end{proposition}

\begin{theorem}[\lean{Consolidation.d\_s\_iff\_exact}]
For a unit $s$ and an element $\sg$,
\[
  \bigl(\forall i,\ \partial_i s = s\cdot \operatorname{sig}\sg_i\bigr)
  \iff
  \bigl(\forall i,\ \operatorname{logDeriv}(s)_i = \operatorname{sig}\sg_i\bigr).
\]
\end{theorem}

A2's derivation condition holds exactly when the unit's log-derivative is
\emph{exact}. So A2 is a unit plus a potential; \Afour{} (one unit per direction)
already supplies the unit; and no geometric result needs the potential
(\lean{Amendment.same\_unit\_same\_geometry}).

The consolidated list is
\[
  \mathrm{A1}\ \cdot\ \mathrm{A2}^{*}\ (\text{one family of units})\ \cdot\
  \mathrm{A5}\ \cdot\ \mathrm{A6}\ \cdot\ \App\ \cdot\ \mathrm{A7},
\]
recorded as \lean{Consolidation.consolidation}. The reduction is not claimed to be
\emph{forced}: a framework may keep a potential primitive. What has changed is
that the potential is now known to be the only surplus.

\section{The source of the geometry cannot be a topological charge}
\label{sec:topology}

\App{} reads $\Delta\sg = \Delta\cdot\nu$ and was presented as an \emph{index}
law, with $\nu$ a topological charge. The development tabulated three candidate
integers. All three are eliminated by one theorem.

\begin{theorem}[\lean{DegreeSource.additive\_nonneg\_on\_group\_is\_trivial}]
Let $G$ be a group and let $f\colon G\to\mathbb{R}$ satisfy $f(ab)=f(a)+f(b)$ and
$f\geq 0$. Then $f\equiv 0$.
\end{theorem}

\noindent The proof is three lines: $f(1)=0$, then $f(a)+f(a^{-1})=0$ with both
terms non-negative.

Additivity is what makes a source an index (\lean{Flux.flux\_combine});
non-negativity is what universal attraction requires
(\lean{Attraction.attraction\_from\_counting}). Hence:

\begin{claimbox}
No group-valued charge can be a universally attractive additive source. Every
topological charge here is group-valued --- $\pi_1$ and $\pi_2$ are groups --- so
the source of the geometry cannot be topological at all.
\end{claimbox}

The argument is sharp rather than over-strong.
\lean{DegreeSource.monoid\_source\_exists} exhibits an additive, non-negative,
non-zero map on $\mathbb{N}$: it is \emph{inverses} that do the killing. What
survives is a threshold count, a cardinality with no anti-content. Two external
checks: real gravity's source is not topological and mass is non-negative; and
electric charge \emph{is} signed and \emph{does} source an interaction, which the
theorem permits because that interaction is not universally attractive.

\begin{remark}[Consequence for \App{}]
The coefficient in $\kappa=-2(n-1)\Delta$ is algebraic in $\nu$ and is untouched.
What is touched is \App{}'s standing: its authority came from $\nu$ being a
topological index. \App{} keeps its consequences and loses its index reading; it
is a postulate of the same standing as the rest of the axioms.
\end{remark}

\section{The pairing was never chosen: two metrics}
\label{sec:twometrics}

The framework computes curvature in two places.

\begin{itemize}[leftmargin=1.4em]
\item \lean{Conformal.RscBare} $=\sum_b \Ric_{bb}$, described in its own
      documentation as the \emph{$\delta$-trace} of Ricci, on the metric
      $e^{2\sg}\delta$. This is where $\kappa=-2(n-1)\Delta$ is derived, so
      \textbf{$\kappa$ is a Euclidean result}.
\item \lean{Diagonal.Ric} on the metric $\mathrm{met}\,\eta\,w$, with the
      signature carried as a parameter satisfying $\eta^{2}=1$. The development's
      own chain table runs the \emph{entire} gravity chain through this sector,
      and the key cancellation requires $\eta_t\eta_r=-1$.
\end{itemize}

The single theorem joining the two, \lean{Anisotropic.ChrDir\_of\_isotropic},
scopes itself in its own documentation to \emph{Euclidean signature}, with
hypothesis $\eta\equiv 1$.

\begin{proposition}[\lean{TwoMetrics.euclidean\_excludes\_lorentzian}]
If $2\neq 0$ in the ring, then $\eta\equiv 1$ and $\eta_t\eta_r=-1$ cannot both
hold.
\end{proposition}

\begin{claimbox}
The only bridge between the sectors holds exactly where the physics does not.
$\kappa$ and the observables have never been shown to be about the same geometry.
\end{claimbox}

\lean{TwoMetrics.delta\_trace\_ne\_eta\_trace} shows the difference does not
cancel: on a Lorentzian pair the two traces read $2$ and $0$.

\begin{remark}[A correction to the record]
The framework's summary of what rests on the axioms alone listed ``the Lorentzian
signature''. What is derived is a \emph{codimension-one splitting} of the
directions from the drift. A splitting is a flag; a signature is a property of a
form. That file's own closing paragraph already granted that the timelike sign is
not derived.
\end{remark}

\section{The repair: the coefficient survives, the operator does not}
\label{sec:etatrace}

The curvature identity was redone from the connection upward with a signature,
since inserting $\eta$ into the final trace alone would not be the same geometry.
The connection used is the one the development already had for the anisotropic
sector, at the isotropic locus with $\eta$ left general: a case that existed and
had never been computed.

\begin{theorem}[\lean{EtaTrace.RscE\_eq}]
\[
  \Rsc^{\eta} \;=\; -2(n-1)\,\Box\sg \;-\; (n-1)(n-2)\,|\nabla\sg|^{2}_{\eta},
\]
where $\Box\sg=\sum_i \eta_i\sg_{ii}$ and
$|\nabla\sg|^{2}_{\eta}=\sum_i \eta_i\sg_i\sg_i$.
\end{theorem}

\emph{Identical coefficients.} \lean{EtaTrace.A6'\_from\_index\_eta} then returns
$\kappa=-2(n-1)\Delta$.

The reason is structural rather than fortunate. In the two-connection contraction
two $\eta$-carrying terms cancel identically, the second carrying $\eta_e^2=1$;
every surviving $\eta$ is either squared or pinned by a $\delta$ that forces its
indices equal. The signature cannot reach the coefficient.

\subsection{But the operator changes}

\begin{theorem}[\lean{Hyperbolic.lapE\_lorentzian\_split}]
With $\eta_t=-1$ and $\eta_i=1$ for $i\neq t$,
$\ \Box\sg=\sum_{i\neq t}\sg_{ii}-\sg_{tt}$.
\end{theorem}

\begin{claimbox}
$\kappa$ is safe; the operator it multiplies is the d'Alembertian. The source law
$\Box\sg = \Delta\cdot\nu$ is \textbf{hyperbolic}, not elliptic --- a sourced wave
equation. Newton is recovered as its static limit, where $\sg_{tt}=0$ and the
d'Alembertian is the Laplacian, and the framework's gravity results are all
computed under staticity.
\end{claimbox}

\lean{Hyperbolic.lap\_sub\_lapE} measures the gap as $\sum_i(1-\eta_i)\sg_{ii}$,
zero only for a Euclidean signature. Consequently the framework's Poisson-limit
theorem should be read as \emph{scoped to static configurations}. Nothing computed
from it changes; what changes is what it is a limit of.

\section{Null characteristics, and what a scale wave is}
\label{sec:waves}

A hyperbolic law has characteristics. Since no integration theory is available,
the wave is treated algebraically: a \emph{wave profile} is a scale whose Hessian
is rank one along a covector $k$, that is $\sg_{ij}=k_ik_j m$.

\begin{theorem}[\lean{Hyperbolic.lapE\_of\_waveProfile},
\lean{Hyperbolic.vacuum\_wave\_is\_null}]
$\Box\sg=\operatorname{bareForm}(\eta,k)\cdot m$, and in vacuum with $m$ a unit,
$\operatorname{bareForm}(\eta,k)=0$.
\end{theorem}

That bare form is the framework's \emph{own} quadratic form, defined kinematically
as the locus where the measured form vanishes. Nothing in the development had said
that anything \emph{travels} along it.

\begin{claimbox}
Scale disturbances propagate on the framework's own null cone: the dynamical
characteristic cone and the kinematic null cone coincide.
\end{claimbox}

\subsection{The effect is quantised by the threshold spectrum}

Content in SCD is what a probe at a given scale resolves, namely
$\mathrm{resolved}(\mu,t)=\{k\mid \mu_k\leq t\}$.

\begin{theorem}[\lean{Directional.resolved\_eq\_iff\_no\_threshold\_between}]
For $t\leq t'$, $\ \mathrm{resolved}(\mu,t)=\mathrm{resolved}(\mu,t')$ if and only
if no $\mu_k$ lies in $(t,t']$.
\end{theorem}

\begin{claimbox}
A scale wave has no continuous observable. It acts only where it crosses a
threshold, so half the amplitude gives \emph{none} of the effect unless a
threshold lies in the interval. A scale wave is a wave of emergence.
\end{claimbox}

\subsection{The breathing mode: expressible, but not radiable}

The development contained a theorem excluding a breathing (scalar) polarization,
justified by the claim that a uniform rescaling is a fiducial shift and so
unobservable by A5.

\begin{proposition}[\lean{Breathing.nonconstant\_scale\_is\_expressible},
\lean{Breathing.uniform\_across\_directions\_still\_varies}]
Fiducial invariance quantifies over one constant added at \emph{every point}, and
the difference pattern compares \emph{points}. A breathing wave is uniform across
\emph{directions} and varies across \emph{points}, and is separated by a
fiducial-invariant functional.
\end{proposition}

So A5 does not remove the breathing mode, and the null-cone result exhibits it.
The conclusion nevertheless survives, on a ground the framework already had: a
scalar disturbance is sourced by the monopole; a conserved source's monopole is
frozen (\lean{Waves.monopole\_does\_not\_radiate}); and a non-conserved source
admits \emph{no field equation at all}
(\lean{Dynamics.no\_field\_equation\_of\_nonconserved}).

\begin{claimbox}
The breathing mode is expressible but cannot be radiated by any source for which a
field equation can be written. Same prediction, different ground, and a narrower
scope: a free or primordial breathing wave is excluded by nothing here.
\end{claimbox}

\begin{remark}[Withdrawn]
An earlier statement of this pass --- \emph{no interferometer can see a scale
wave} --- is withdrawn. The cited theorem takes a \emph{single} scale pattern and
compares two null directions at one point; it excludes finding a \emph{static}
anisotropy. A wave differs between the arms, and between emission and return. What
survives is that a static scale pattern is invisible to an interferometer at any
precision.
\end{remark}

\section{Off the commuting locus: an antisymmetric Ricci without torsion}
\label{sec:quantum}

Every quantitative result in the development is computed over a commutative ring,
where the adjoint action vanishes identically and the quantum correction is zero
by construction. This section leaves that locus.

The non-commutative Riemann tensor is known: curvature is still the
Kulkarni--Nomizu product of the bare bookkeeping with the scale deformation, plus
one correction built from commutators of scale gradients. Contracting it:

\begin{proposition}[\lean{QuantumRicci.sum\_RmCorr\_contract}]
The Riemann correction vanishes in the Ricci contraction: its two Kronecker terms
cancel term by term at $c=a$.
\end{proposition}

\begin{theorem}[\lean{QuantumRicci.two\_Ric\_eq\_nc}]
$2\,\Ric_{be} = (2-n)\,\Defm_{be} - \delta_{be}\operatorname{tr}\Defm$, with no
extra term.
\end{theorem}

Yet the deformation tensor is itself not symmetric off the locus, and therefore:

\begin{theorem}[\lean{QuantumRicci.two\_Ric\_antisymm}]
\[
  \boxed{\ \Ric_{be}-\Ric_{eb} \;=\; (n-2)\,[\sg_b,\sg_e]\ }
\]
\end{theorem}

\begin{theorem}[\lean{QuantumRicci.two\_RscBare\_eq\_nc}]
$2\,\Rsc = (2-2n)\operatorname{tr}\Defm$, and $\operatorname{tr}\Defm$ is built
from $\sg_{aa}$, $\sg_a\sg_a$ and $|\nabla\sg|^{2}$ --- every term a square.
\end{theorem}

\begin{claimbox}
The entire quantum correction to the geometry is the antisymmetric part of Ricci.
The scalar curvature does not move at all.
\end{claimbox}

\subsection{What this predicts}

\paragraph{An antisymmetric Ricci with a torsion-free connection.}
The connection is symmetric in its lower pair whether or not the values commute
(\lean{Conformal.Chr\_symm}, proved over a general ring), so there is \emph{no
torsion}. In ordinary differential geometry $\Ric_{[ab]}=0$ for a torsion-free
connection, by the first Bianchi identity --- and that proof uses commutativity of
the \emph{values}, not only torsion-freeness. This is therefore not
Einstein--Cartan, which buys the same antisymmetry by giving the connection
torsion.

\paragraph{Nothing classical moves, at any order.}
The source law, $\kappa$, the Newtonian limit and every cosmological statement are
exactly their commutative selves. In particular there is \textbf{no $\hbar G/r^{3}$
correction to the Newtonian potential}, of the kind effective-field-theory
treatments of gravity produce, and \textbf{no running of $G$}. This is a flat
disagreement with the standard expectation, and it is exact rather than small.

\paragraph{It vanishes in two dimensions,}
by the same factor $(n-2)$ that governs conformal geometry throughout.

\paragraph{Where it can appear.}
An antisymmetric Ricci in a field equation requires an antisymmetric source, and
the physical antisymmetric source is spin. Independently,
\lean{NCConformal.quantum\_correction\_is\_a\_wedge} holds by \lean{rfl}: the
correction \emph{is} the framework's rotational label applied to the scale
gradient pattern. \emph{This is an identification, not a theorem}: no field
equation for the antisymmetric half is written, and no magnitude is claimed.

\paragraph{And it does not inherit the pairing problem.}
By \lean{Micro.quantum\_correction\_signature\_independent} the signature enters
the deformation tensor only symmetrically, so the antisymmetric part --- where the
whole correction lives --- is the same for every $\eta$. Unlike $\kappa$, this
prediction never depended on the pairing that was never chosen.

\section{Constants and dimensions}
\label{sec:constants}

The usual bookkeeping has three base dimensions $L,T,M$ and four constants
$c,\hbar,G,e$. Here there are no dimensions at all: what replaces dimension is a
single $\mathbb{Z}$-grading, the weight under $\sg\mapsto c\sg$.

The collapse is structural rather than a choice of units. Time is a
\emph{direction}, told from space by a sign and not by a dimension. Light crosses
at measured speed one in every direction, because the scale sets rods and light
times alike, so $c=1$ is a theorem. And A3, being a group law
(\S\ref{sec:consolidation}), places mass on the same axis with the opposite sign.

The magnitudes on that axis are $\Delta$ (weight $1$), the scale step identified
with $\hbar$ (weight $1$, since a step is a difference of log-scales), and the
threshold density $\rho$ (weight $-1$).

\begin{theorem}[\lean{Constants.two\_pure\_numbers},
\lean{Constants.collapse\_iff}]
Both $\Delta\rho$ and $\hbar\rho$ are weight zero, and they coincide if and only
if $\Delta=\hbar$.
\end{theorem}

\begin{claimbox}
There are \emph{two} free pure numbers, not one, unless $G$ and $\hbar$ are the
same magnitude --- and nothing in the development says they are.
\end{claimbox}

The framework's own files already record that ``$G$ and $\hbar$ share an origin''
is conditional on a step not taken, and that ``$\hbar$ is that step'' is an
identification rather than a construction. The two statements had never been read
against each other.

\subsection{\texorpdfstring{Where $e$ goes}{Where e goes}}

The fine-structure constant is weight zero: a gauge coupling's value does not
depend on how $\sg$ is labelled. So it is a pure number \emph{structurally}. And
the asymmetry has a reason: the gauge law relates weight-zero quantities to
weight-zero quantities, while gravity's relates a weight-zero \emph{count} to a
weight-one \emph{scale}.

\begin{claimbox}
Exactly one law crosses the grading. ``Why is $\alpha$ dimensionless and $G$ not''
stops being a fact about dimensions and becomes a fact about which law crosses the
grading.
\end{claimbox}

\begin{table}[t]
\centering
\begin{tabular}{@{}ll@{}}
\toprule
\textbf{Usual} & \textbf{Here} \\
\midrule
$L,T,M$ & one grading: weight under $\sg\mapsto c\sg$ \\
$c$ & not a constant --- one cone; measured speed one is a theorem \\
$\hbar$ & a weight-one scale step; an identification, not a construction \\
$G$ & a weight-one coefficient $\Delta$ \\
$e$, i.e.\ $\alpha$ & weight zero --- already pure, and structurally so \\
Planck units & one axis, so one magnitude fixes everything \\
pure numbers & $\Delta\rho$ and $\hbar\rho$: \emph{two}, collapsing iff
$\Delta=\hbar$ \\
\bottomrule
\end{tabular}
\caption{The dimensional bookkeeping, translated.}
\end{table}

\begin{remark}[A question the usual picture cannot ask]
$[G]=L^{3}M^{-1}T^{-2}$ and $[\hbar]=ML^{2}T^{-1}$ have different dimensions, so
\emph{are $G$ and $\hbar$ the same magnitude?} is not a well-formed question there.
Here they carry the same weight, the question is well-formed, and its answer is a
single pure number.
\end{remark}

\section{What remains open}
\label{sec:open}

\paragraph{The directional source law.}
\Ap{} reads $e^{2\sg}\Rsc=\kappa\rho$, and $e^{2\sg}$ is \emph{one} conformal
factor. With a scale per direction there is no single factor:
\lean{Directional.dirTrace\_factors\_iff} shows the factor comes out only when
every directional unit agrees. So \Ap{}'s \emph{form} presupposes an isotropic
scale while \Afour{} supplies a directional one. A replacement is written
(\lean{Breathing.DirSource}) on the true scalar curvature --- \Afour{} supplies
exactly the invertibility the bare form was contorted to avoid --- but the
curvature identity that would let it produce a number does not exist off the
isotropic locus. Until it does, $\kappa$ and the observables meet only where the
scale is isotropic.

\paragraph{The quantum correction has never been carried.}
Section~\ref{sec:quantum} is the first computation off the commuting locus. Every
number in the development is computed on the isotropic, commuting locus, and both
unfinished sectors are off it.

\paragraph{\texorpdfstring{$\Delta=\hbar$}{Delta = hbar}: deliberately left open.}
Nothing in the axioms relates the source coefficient to a scale shift, and the
scale shift itself is not built from the axioms --- it requires a ring
endomorphism and a unit that nothing supplies. The identification would have to
come from outside. It is recorded as an open question, and the free-number count
reads \emph{two}.

\paragraph{And one number that cannot be computed.}
Predictions are informationally disjoint from inputs, and counts cannot fix a
unit; $\Delta\rho$ is a \emph{permanent} input. The framework therefore predicts
ratios and structures, not magnitudes, and the familiar arcsecond figures are
weight-zero statements evaluated on measured inputs.

\section{Method}

Three times in this pass a result was kept while its stated ground was replaced:
\App{} keeps its consequences and loses its index reading; the Poisson limit keeps
its numbers and becomes a static limit; the breathing exclusion keeps its
conclusion and changes its argument. Twice a file's own self-correction caught an
error that no theorem would have. Once a claim made in this pass was withdrawn one
section later.

The rate is the thing to read, not the instances. A formalisation makes the
theorems reliable; it does not make the \emph{prose around them} reliable, and
most of what was corrected here lived in the prose --- in what a true theorem was
being quoted as showing.

\end{document}
